\documentclass[11pt]{article}

\usepackage[margin=1in]{geometry}
\usepackage{booktabs}
\usepackage{makecell}
\usepackage{graphicx}
\usepackage{amsmath}
\usepackage{caption}
\usepackage{float}
\usepackage{siunitx}
\usepackage[hidelinks]{hyperref}

\graphicspath{{img/}}
\captionsetup{font=small}

\title{\textbf{A Two-Language Replication of}\\
Acemoglu, Johnson, and Robinson (2001),\\
\emph{``The Colonial Origins of Comparative Development''}}
\author{Replication report}
\date{\today}

\begin{document}
\maketitle

\begin{abstract}
\noindent We replicate the core empirical results of Acemoglu, Johnson, and
Robinson (2001, \emph{AER}) --- the three scatter relationships (Figures 1--3)
and the descriptive, OLS, and instrumental-variables tables (Tables 1, 2, and
4). The replication is carried out \emph{twice}, independently, in Python and in
R, reading a single shared copy of the authors' data. Both implementations
recover the paper's headline two-stage least-squares (2SLS) estimate of the
effect of institutions on income, $0.94$ (s.e.\ $0.16$), and agree with each
other to the last reported digit. The one substantive cross-language difference
we encountered --- the standard error of the 2SLS coefficient --- turned out to
be a degrees-of-freedom convention, which we reconcile so that both languages
match the printed paper exactly.
\end{abstract}

\section{The paper in one paragraph}

AJR ask whether \emph{institutions} cause long-run prosperity. The empirical
problem is reverse causality and omitted variables: rich countries can afford
good institutions, so a raw correlation between institutional quality and income
is not evidence of a causal effect. AJR's identification strategy uses
\emph{settler mortality} as an instrument for institutions. Where European
settlers faced high mortality (tropical disease environments) they set up
``extractive'' states rather than settling; where mortality was low they
reproduced European-style institutions protecting property rights. Settler
mortality plausibly affects income today \emph{only} through its historical
effect on institutions, making it a valid instrument. The causal chain is:
\[
\text{settler mortality} \;\longrightarrow\; \text{early institutions}
\;\longrightarrow\; \text{current institutions} \;\longrightarrow\;
\text{income today}.
\]
The three key variables are \texttt{logpgp95} (log GDP per capita, PPP, 1995),
\texttt{avexpr} (average protection against expropriation risk, 1985--95, the
institutions proxy), and \texttt{logem4} (log settler mortality, the
instrument). The main analysis runs on a ``base sample'' of 64 former colonies.

\section{What we replicated}

\begin{itemize}
  \item \textbf{Figure 1} --- reduced form: income vs.\ settler mortality.
  \item \textbf{Figure 2} --- OLS: income vs.\ expropriation risk.
  \item \textbf{Figure 3} --- first stage: expropriation risk vs.\ settler mortality.
  \item \textbf{Table 1} --- descriptive statistics by sample and mortality quartile.
  \item \textbf{Table 2} --- OLS regressions of income on institutions.
  \item \textbf{Table 4} --- 2SLS regressions (the headline result), with first
        stage and OLS panels.
\end{itemize}

\section{Figures}

The three figures reproduce the paper's bivariate relationships on the base
sample ($n=64$), plotting country codes as markers with an OLS fit line. The
fitted slopes match the corresponding regression coefficients in the tables: the
Figure~2 slope ($0.522$) is the Table~2 column~(2) OLS coefficient, and the
Figure~3 slope ($-0.607$) is the Table~4 first-stage coefficient.

\begin{figure}[H]
  \centering
  \includegraphics[width=0.82\textwidth]{figure1.png}
  \caption{Reduced-form relationship between income and settler mortality
  (slope $-0.573$, s.e.\ $0.076$, $R^2=0.477$). High-mortality colonies are
  poorer today.}
\end{figure}

\begin{figure}[H]
  \centering
  \includegraphics[width=0.82\textwidth]{figure2.png}
  \caption{OLS relationship between expropriation risk and income
  (slope $0.522$, s.e.\ $0.061$, $R^2=0.540$). Better property-rights
  protection is associated with higher income.}
\end{figure}

\begin{figure}[H]
  \centering
  \includegraphics[width=0.82\textwidth]{figure3.png}
  \caption{First-stage relationship between settler mortality and expropriation
  risk (slope $-0.607$, s.e.\ $0.127$, $R^2=0.270$). The instrument is strongly
  correlated with the endogenous regressor.}
\end{figure}

\section{Tables}

All tables report the reproduced values; standard errors are in parentheses.
Each matches the printed paper up to rounding. The notes below each table flag
the few known data-driven discrepancies.

\begin{table}[H]
  \centering
  \caption{Descriptive statistics (mean; standard deviation in parentheses).
  Columns (1)--(4) split the base sample by quartiles of settler mortality.}
  \footnotesize
  \resizebox{\textwidth}{!}{%
  \begin{tabular}{lcccccc}
    \toprule
    & Whole world & Base sample & (1) & (2) & (3) & (4) \\
    \midrule
    Log GDP per capita (PPP) 1995              & 8.30 (1.11) & 8.06 (1.04) & 8.89 (1.29) & 8.41 (0.64) & 7.78 (0.80) & 7.20 (0.63) \\
    Log output per worker 1988                 & -1.73 (1.08) & -1.93 (0.98) & -1.03 (0.86) & -1.46 (0.42) & -2.20 (0.80) & -3.03 (0.48) \\
    Avg.\ expropriation risk 1985--95          & 7.07 (1.80) & 6.52 (1.47) & 7.90 (1.55) & 6.47 (0.95) & 5.97 (1.32) & 5.89 (1.30) \\
    Constraint on executive 1990               & 3.64 (2.34) & 3.97 (2.26) & 5.33 (2.10) & 5.12 (1.80) & 3.31 (2.18) & 2.27 (1.58) \\
    Constraint on executive 1900               & 1.86 (1.82) & 2.25 (2.11) & 3.67 (2.99) & 3.28 (2.11) & 1.12 (0.50) & 1.00 (0.00) \\
    Constraint at independence                 & 3.59 (2.41) & 3.40 (2.39) & 5.17 (2.76) & 2.44 (1.92) & 3.12 (2.28) & 3.43 (2.14) \\
    Democracy 1900                             & 1.15 (2.58) & 1.64 (3.00) & 3.92 (4.58) & 2.76 (2.91) & 0.19 (0.40) & 0.00 (0.00) \\
    European settlements 1900                  & 0.30 (0.42) & 0.16 (0.26) & 0.32 (0.44) & 0.26 (0.18) & 0.08 (0.12) & 0.01 (0.02) \\
    Log settler mortality                      & n.a.\ & 4.66 (1.26) & 3.01 (0.59) & 4.28 (0.05) & 4.92 (0.40) & 6.35 (0.75) \\
    \midrule
    Number of observations                     & 163 & 64 & 14 & 18 & 17 & 15 \\
    \bottomrule
  \end{tabular}}
\end{table}

\begin{table}[H]
  \centering
  \caption{OLS regressions. Dependent variable is log GDP per capita 1995, except
  columns (7)--(8) which use log output per worker 1988. ``base'' = 64-country
  base sample; ``world'' = whole-world sample.}
  \footnotesize
  \resizebox{\textwidth}{!}{%
  \begin{tabular}{lcccccccc}
    \toprule
    & (1) & (2) & (3) & (4) & (5) & (6) & (7) & (8) \\
    & world & base & world & world & base & base & world & base \\
    \midrule
    Avg.\ expropriation risk & 0.53 (0.04) & 0.52 (0.06) & 0.46 (0.06) & 0.39 (0.05) & 0.47 (0.06) & 0.40 (0.06) & 0.45 (0.04) & 0.46 (0.06) \\
    Latitude                 &             &             & 0.87 (0.49) & 0.33 (0.45) & 1.58 (0.71) & 0.88 (0.63) &             &             \\
    Asia dummy               &             &             &             & -0.15 (0.15) &            & -0.58 (0.23) &            &             \\
    Africa dummy             &             &             &             & -0.92 (0.17) &            & -0.88 (0.17) &            &             \\
    ``Other'' dummy          &             &             &             & 0.30 (0.37) &             & 0.11 (0.38) &             &             \\
    \midrule
    $R^2$                    & 0.61 & 0.54 & 0.62 & 0.72 & 0.57 & 0.71 & 0.55 & 0.49 \\
    Observations             & 111 & 64 & 111 & 111 & 64 & 64 & 108 & 61 \\
    \bottomrule
  \end{tabular}}
  \caption*{\footnotesize Note: the whole-world $N=111$, not the printed $110$;
  AJR's replication code notes one observation was mistakenly dropped from the
  printed table, so $111$ is the corrected figure (matching standard
  reproductions).}
\end{table}

\setcounter{table}{3}% this is AJR's Table 4 (we replicate AJR Tables 1, 2, 4)
\begin{table}[H]
  \centering
  \caption{IV regressions of log GDP per capita --- the headline result. Average
  expropriation risk is instrumented by log settler mortality. Column (9) uses
  log output per worker. America is the omitted continent dummy.}
  \footnotesize
  \resizebox{\textwidth}{!}{%
  \begin{tabular}{lccccccccc}
    \toprule
    & (1) & (2) & (3) & (4) & (5) & (6) & (7) & (8) & (9) \\
    & base & base & \makecell{no Neo-\\Europes} & \makecell{no Neo-\\Europes} & \makecell{no\\Africa} & \makecell{no\\Africa} & \makecell{+\,cont.\\dummies} & \makecell{+\,cont.\\dummies} & \makecell{output/\\worker} \\
    \midrule
    \multicolumn{10}{l}{\emph{Panel A: Two-Stage Least Squares}} \\
    Avg.\ expropriation risk & 0.94 (0.16) & 1.00 (0.22) & 1.28 (0.36) & 1.21 (0.35) & 0.58 (0.10) & 0.58 (0.12) & 0.98 (0.30) & 1.11 (0.46) & 0.98 (0.17) \\
    Latitude                 &             & -0.65 (1.34) &            & 0.94 (1.46) &            & 0.04 (0.84) &            & -1.18 (1.76) &            \\
    Asia dummy               &             &             &             &             &            &             & -0.92 (0.40) & -1.05 (0.52) &           \\
    Africa dummy             &             &             &             &             &            &             & -0.46 (0.36) & -0.44 (0.42) &           \\
    ``Other'' dummy          &             &             &             &             &            &             & -0.94 (0.85) & -0.99 (1.00) &           \\
    \midrule
    \multicolumn{10}{l}{\emph{Panel B: First Stage (dependent variable: expropriation risk)}} \\
    Log settler mortality    & -0.61 (0.13) & -0.51 (0.14) & -0.39 (0.13) & -0.39 (0.14) & -1.21 (0.22) & -1.14 (0.24) & -0.43 (0.17) & -0.34 (0.18) & -0.63 (0.13) \\
    $R^2$                    & 0.27 & 0.30 & 0.13 & 0.13 & 0.47 & 0.47 & 0.30 & 0.33 & 0.29 \\
    \midrule
    \multicolumn{10}{l}{\emph{Panel C: Ordinary Least Squares}} \\
    Avg.\ expropriation risk & 0.52 (0.06) & 0.47 (0.06) & 0.49 (0.08) & 0.47 (0.07) & 0.48 (0.07) & 0.47 (0.07) & 0.42 (0.06) & 0.40 (0.06) & 0.46 (0.06) \\
    \midrule
    Observations             & 64 & 64 & 60 & 60 & 37 & 37 & 64 & 64 & 61 \\
    \bottomrule
  \end{tabular}}
  \caption*{\footnotesize Note: the public \texttt{maketable4.dta} does not ship
  the ``Other'' continent dummy used in columns (7)--(8); we merge it in from
  \texttt{maketable2.dta} on the country code. Panel A standard errors use the
  small-sample ($n-k$) correction, matching the printed paper.}
\end{table}

\section{Reading the results}

The headline comparison is OLS vs.\ IV. In Table~2 the OLS coefficient on
expropriation risk is about $0.52$ (base sample); in Table~4 Panel~A the 2SLS
coefficient is $0.94$. The instrumented estimate is \emph{larger}, not smaller.
AJR read this as evidence that OLS is, if anything, attenuated by measurement
error in the institutions proxy, and that the cross-country correlation is not
merely reverse causation. Panel~B confirms a strong first stage --- log settler
mortality enters with a coefficient of $-0.61$ (s.e.\ $0.13$) in the base
sample --- so the instrument is not weak. The result is robust across samples:
dropping the ``Neo-Europes'' (columns 3--4) or all of Africa (columns 5--6), or
adding continent dummies (columns 7--8), leaves a large, significant effect.

\section{The two-language replication}

We implemented the entire pipeline twice, independently, against one shared copy
of the data (\texttt{data/maketable\{1,2,4\}.dta}).

\begin{table}[H]
  \centering
  \caption*{\textbf{Tools used in each implementation.}}
  \small
  \begin{tabular}{lll}
    \toprule
    Step & Python & R \\
    \midrule
    Read Stata \texttt{.dta} & \texttt{pandas.read\_stata} & \texttt{foreign::read.dta} \\
    Figures                  & \texttt{matplotlib}        & base graphics \\
    OLS (Tables 1--2)        & \texttt{statsmodels}       & \texttt{stats::lm} \\
    2SLS (Table 4)           & \texttt{linearmodels.IV2SLS} & \texttt{ivreg::ivreg} \\
    \bottomrule
  \end{tabular}
\end{table}

\noindent The two stacks were chosen to be idiomatic rather than identical. On
the R side, figures and the OLS tables need only base R plus the bundled
\texttt{foreign} package --- no installation --- while Table~4 adds the
lightweight \texttt{ivreg} package.

\paragraph{Agreement.} Every point estimate agrees between Python and R to the
two decimal places reported, and both match the paper. The figures' slopes,
$R^2$, and sample sizes are identical across languages.

\paragraph{The one wrinkle: 2SLS standard errors.} Our first R and Python
versions disagreed on a single quantity --- the standard error of the headline
2SLS coefficient: R reported $0.16$, Python $0.15$. This was not a bug but a
\emph{degrees-of-freedom convention}. R's \texttt{ivreg} divides the residual
variance by $n-k$ (the textbook small-sample correction AJR used), whereas
\texttt{linearmodels} with \texttt{cov\_type="unadjusted"} divides by $n$.
Setting \texttt{debiased=True} in the Python fit applies the $n-k$ correction, after
which both languages report $0.94\ (0.16)$ --- matching the printed paper
exactly. This is a useful illustration of why a second, independent
implementation is valuable: it surfaced a silent modeling assumption that a
single-language replication would never have questioned.

\section{Data notes and minor discrepancies}

\begin{itemize}
  \item \textbf{Whole-world $N=111$, not $110$.} AJR's own replication code notes
        one observation was mistakenly dropped from the printed Table~2; $111$
        is the corrected sample, and matches standard reproductions.
  \item \textbf{The ``Other'' continent dummy} is absent from the public
        \texttt{maketable4.dta}, so we merge it from \texttt{maketable2.dta} on
        the country code for Table~4 columns (7)--(8).
  \item \textbf{Mortality-quartile boundaries (Table 1).} Five countries are
        stored at the float value of $78.1$; the paper assigns them to quartile
        (2), so the (2)/(3) cut is placed at $78.15$ to reproduce the
        $14/18/17/15$ split.
\end{itemize}

\section{Reproducibility}

From the \texttt{paper-replication/} directory:
\begin{quote}
\texttt{make python}\quad --- run all Python scripts ($\rightarrow$ \texttt{python/figures}, \texttt{python/tables})\\
\texttt{make R}\qquad\;\, --- run all R scripts ($\rightarrow$ \texttt{R/figures}, \texttt{R/tables})\\
\texttt{make all}\qquad --- both
\end{quote}
Each language writes its figures as \texttt{.png} and its tables as plain-text
files; the values in this report are taken directly from those outputs.

\bigskip
\noindent\emph{Reference:} Acemoglu, D., Johnson, S., \& Robinson, J.~A. (2001).
The Colonial Origins of Comparative Development: An Empirical Investigation.
\emph{American Economic Review}, 91(5), 1369--1401.

\end{document}
